\documentclass[noindent,nohyp,parspace,titlepage,twoside,12pt]{article}

\usepackage{algorithm}
\usepackage{algpseudocode}
\usepackage{amsfonts}
\usepackage{amsmath}
\usepackage[backend=biber,sorting=none]{biblatex}
% \usepackage[type={CC},modifier={by-sa},version={4.0}]{doclicense}
\usepackage{geometry}
\usepackage{hyperref}
\usepackage{pgfplots}
\usepackage{xparse}

\addbibresource{thesis.bib}

\usetikzlibrary{shapes,matrix}

\hypersetup{
  colorlinks,
  citecolor=black,
  filecolor=black,
  linkcolor=black,
  urlcolor=black
}

\pgfplotsset{compat=1.18}
\setcounter{secnumdepth}{5}
\setcounter{tocdepth}{5}

\def\TITLE{\
  Investigating Bias \\
  in LLM Self-Evaluation\
}
\def\AUTHOR{Attila M. Magyar}

\title{\TITLE}
\author{\AUTHOR}

\linespread{1.25}
\begin{document}

\begin{titlepage}

  \begin{center}
    \Huge\textbf{\TITLE}\normalsize
  \end{center}
  \begin{center}
    \Large Thesis \normalsize
  \end{center}

  \vfill

% TODO: author, supervisor

  \vfill

  \begin{center}
    \Large Mathematics Expert in Data Analytics and Machine Learning
    \normalsize
  \end{center}

  \begin{center}
    \includegraphics[scale=0.2]{img/logo.png} \\
  \end{center}

  \begin{center}
    \Large Eötvös Loránd University \\
    \Large Faculty of Science \normalsize
  \end{center}

  \begin{center}
    \Large Budapest, 2025 \normalsize
  \end{center}

\end{titlepage}

  \tableofcontents

\newpage

  \section{Introduction}

    \subsection{A Brief Introduction to LLMs}

      A language model is a machine learning model designed to perform a wide
      range of tasks that involve natural language processing (NLP), including
      text summarization, translation, sentiment analysis, spam detection,
      content moderation, text generation, etc.

      Significant advancements in deep learning \cite{attention,gpt3,gpt4} led
      to the emergence of \textbf{large language models} (LLMs) --- particularly
      generative LLMs --- which in the early 2020s became productized and
      widely adopted in both industry and popular discourse.

      A generative large language model is a model which has a parameter count
      on the order of billions or more, and predicts the conditional
      probability \cite{llms}

      \begin{align} \label{eqautoreg}
        P(w_m | w_0, \cdots, w_{m-1})
      \end{align}

      where $m \in \mathbb{N}$, $w_0$ is a special start symbol, and $w_k$ is
      the $k$-th token (for $1 \le k \le m$) in a sequence of tokens that form
      a piece of text in some (natural) language. The interpretation of the
      tokens depends on the exact tokenization strategy used, which may define
      tokens as words, word pieces, n-grams, or individual characters, and
      spaces, punctuation marks, etc.

      Text generation then is an autoregressive process where given a
      sequence of tokens as a prefix --- known as the \textbf{prompt} --- the
      model estimates the probability distribution of the next token, takes a
      sample from that distribution, appends it to the sequence, and repeats
      the process with the extended sequence until a given stopping condition
      is met.

      The sampling can usually be controlled via numerical parameters; a
      commonly used one is the \textbf{temperature}: the closer it is to 0, the
      more the sampling will lean toward the most probable token --- making the
      algorithm more deterministic ---, while higher values increase the
      randomization, making the generated text feel more \emph{creative} until,
      above a certain threshold, it becomes incoherent and semantically
      meaningless. If $v \in \mathbb{N}$ denotes the number of all possible
      tokens available for the model (vocabulary size), and
      $\mathbf{s} \in \mathbb{R}^v$ is an output vector of the model assigning
      a score to each token as the continuation of a given input, then the
      distribution for the sampling, with respect to the temperature $T \in
      \mathbb{R}$ is calculated via the softmax function:

      \begin{align} \label{eqsoftmax}
        \text{softmax}(\mathbf{s};T)
          = \left[
              \frac{\exp(\frac{s_i}{T})}{\sum_{j=1}^v \exp(\frac{s_j}{T})}
            \right]_{i=1}^v
      \end{align}

      In practice, if $T$ is sufficiently close or exactly equal to $0$, then
      the sampling is usually replaced with the deterministic argmax function
      in order to preserve numerical stability. Non-zero $T$ values control
      the flatness of the distribution, leading to the aforementioned behavior.

      With sufficiently large model complexity and training corpora size and
      diversity, LLMs start to exhibit capabilities which rival that of top
      performer humans in a broad class of problems \cite{gpt3,gpt4}. The
      versatility of the models is often utilized in a setting where the prompt
      is composed of two parts, each consisting of instructions given in
      natural language:

      \begin{itemize}
        \item the \textbf{system prompt} can instruct the model to behave in a
              certain way, for example, to act like a helpful AI assistant,
              an expert in a domain, or to generate its texts in the style of
              fictional 18th-century Caribbean pirates,

        \item and the \textbf{user prompt} which describes a task to be carried
              out by the model, ranging from text translation or summarization
              to solving complex programming problems or pointing out business
              risks in legal documents, and more.
      \end{itemize}

      Generative models with sufficient generalization capabilities can predict
      likely continuations of such prompts with so high accuracy that the
      generated text will often contain an actual solution to the proposed
      problem. This instruction-following paradigm enables models to perform
      \textbf{few-shot learning} \cite{gpt3} or even \textbf{zero-shot learning}
      by interpreting tasks directly from the natural language description,
      based on just a few or zero examples, respectively, without specific
      training or fine-tuning.

      The problem solving performance of LLMs can be improved further by
      prompt engineering techniques like \textbf{chain-of-thought} prompting
      \cite{cot}, where the model is provided with a step-by-step example
      solution to a related problem in the prompt, encouraging it to also
      articulate intermediate reasoning steps.

    \subsection{LLM Evaluators, LLM-as-a-Judge}

\newpage

  \nocite{*}
  \printbibliography[heading=bibintoc]

\end{document}
